\section{Conclusions}\label{sec:conclusions}

\method supports preliminary sanitary sewer network planning by combining geospatial data, interactive graph editing, elevation enrichment, spatial persistence, and graph-based processing in a single web workflow.
The current version already includes a functional web basis, project persistence, an editor, integration with public street and elevation data, and a graph pipeline capable of generating a directed and simplified preliminary network.

The main partial contribution is the transformation of the method into an operable system in which specialists can create projects, visualize data, edit the mesh, and trigger processing.
By organizing geospatial data, interactive editing, graph processing, and persistence in a single flow, \method provides a reproducible basis for case studies and for the assessment of preliminary network alternatives.

The main remaining limitation is the need for systematic experimental validation with real case studies and qualitative analysis of the generated results.
Thus, the next phase of the project should focus less on basic system infrastructure and more on evaluating the quality of the generated solution, refining metrics, and consolidating acceptance criteria for different urban contexts.
The building-detection component strengthens this direction by adding a machine-learning basis for estimating sewage contribution from the built environment, with future refinements focused on instance counting, lot-level classification, and multimodal urban analysis.
